% A Dimensional Ladder Connecting Tensor Powers of a Finite-Field 4-Algebra
% to Real Clifford Algebras Cl(2n) for n = 0..5
%
% Brayden R. Sanders, 7Site LLC
% Journal-ready (internal language stripped per JOURNAL_LANGUAGE_GUIDE.md)
% Target venues:
%   - Linear Algebra and its Applications (Elsevier)
%   - Journal of Algebra (Elsevier)
%   - Algebra Universalis (Springer)
%
% Verification: test_tig_dirac.py, test T15 + T13 + T14
% MSC 2020: 15A66, 17A30, 11T55

\documentclass[11pt,reqno]{amsart}

\usepackage{amsmath, amssymb, amsthm, mathtools}
\usepackage[margin=1in]{geometry}
\usepackage[colorlinks=true,linkcolor=blue,citecolor=blue,urlcolor=blue]{hyperref}
\usepackage{microtype}
\usepackage{booktabs}

\theoremstyle{plain}
\newtheorem{theorem}{Theorem}[section]
\newtheorem{proposition}[theorem]{Proposition}
\newtheorem{lemma}[theorem]{Lemma}

\theoremstyle{definition}
\newtheorem{definition}[theorem]{Definition}

\theoremstyle{remark}
\newtheorem*{remark}{Remark}

\newcommand{\Z}{\mathbb{Z}}
\newcommand{\F}{\mathbb{F}}
\newcommand{\R}{\mathbb{R}}
\newcommand{\Cl}{\mathrm{Cl}}

\title[Dimensional Ladder $V^{\otimes n} \leftrightarrow \Cl(2n)$]
{A Dimensional Ladder Connecting Tensor Powers of a \\
 Finite-Field 4-Algebra to Real Clifford Algebras $\Cl(2n)$}

\author{Brayden R.~Sanders}
\address{7Site LLC, Hot Springs, Arkansas, USA}
\email{brayden@7site.co}

\subjclass[2020]{15A66, 17A30, 11T55}
\keywords{Clifford algebra, tensor power, finite-field algebra,
binomial decomposition, geometric algebra}

\date{\today}

\begin{document}

\begin{abstract}
We exhibit a dimensional ladder relating the tensor powers $V^{\otimes n}$
of a 4-dimensional commutative non-associative algebra $V$ over $\F_5$ to
the real Clifford algebras $\Cl(2n)$:
\[
\dim_{\F_5} V^{\otimes n} = 4^n = 2^{2n} = \dim_\R \Cl(2n)
\quad\text{for } n = 0, 1, \ldots, 5.
\]
The dimensional match is forced by $\dim V = 4 = 2^2$. The non-trivial
content is that the binomial cell decomposition of $V^{\otimes n}$ ---
arising from the idempotent decomposition $V = (\F_5 e_2 \oplus \F_5 e_3)
\oplus (\F_5 e_0 \oplus \F_5 e_4)$ --- corresponds to the grade
decomposition of $\Cl(2n)$.

At $n = 5$ in particular, the binomial $1 + 5 + 10 + 10 + 5 + 1 = 32$
exactly matches the representation content of $\mathbf{1} \oplus
\bar{\mathbf{5}} \oplus \mathbf{10}$ plus its conjugate, the
representations carried by one fermion generation plus its antimatter
in the SU(5) Grand Unified Theory.

The paper presents the algebra and the dimensional ladder; it does not
assert physical interpretations of the SU(5) decomposition. The
dimensional ladder is verified at $n = 0, 1, 2, 3, 4, 5$ in the
accompanying script \texttt{test\_tig\_dirac.py}.
\end{abstract}

\maketitle

\tableofcontents

%==============================================================
\section{Introduction}
%==============================================================

The real Clifford algebras $\Cl(p, q)$ over a quadratic form of signature
$(p, q)$ form a periodic family under the operation
$\Cl(p, q) \otimes \Cl(0, 0) \cong \Cl(p, q)$
and Bott periodicity $\Cl(n + 8) \cong \Cl(n) \otimes \Cl(8)$.  Their
applications span geometric algebra (Hestenes-Sobczyk
\cite{HesteneSobczyk1984}), spinor representations of orthogonal groups,
quantum-information stabilizer codes, and Standard Model gauge structure.

This paper studies a 4-dimensional commutative non-associative algebra
$V$ over $\F_5$ whose tensor powers exhibit dimensions matching $\Cl(2n)$
exactly for $n=0..5$. The dimensional ladder is forced by $\dim V = 4
= 2^2$; we show that the additional binomial cell structure inside
$V^{\otimes n}$ aligns with the grade decomposition of $\Cl(2n)$.

\subsection{Notation.}
Throughout, $V$ denotes the 4-dimensional commutative non-associative
algebra over $\F_5$ studied in our companion paper \cite{Sanders2026Dirac};
the multiplication is given by an explicit $4 \times 4$ table on a
specific basis $\{e_0, e_2, e_3, e_4\}$.  The structure of $V$ is
summarized in Section~\ref{sec:V}; full details are in
\cite{Sanders2026Dirac}.

%==============================================================
\section{The algebra $V$}
\label{sec:V}
%==============================================================

The 4-dimensional algebra $V$ over $\F_5$ has 4 idempotents (one zero, three
non-zero), a 1-dim Grassmann annihilator $\F_5 \cdot e_0$, and a 1-dim
associator image. Its key features for the present paper are:
\begin{itemize}
\item $V = (\F_5 e_2 \oplus \F_5 e_3) \oplus (\F_5 e_0 \oplus \F_5 e_4)$
as a direct sum (as $\F_5$-vector space), where $e_2$ is a non-zero
idempotent.
\item Each tensor factor of $V^{\otimes n}$ thus carries a binary "sign"
$\pm$ indicating which of the two summands it contributes to.
\item Across $n$ tensor factors, this gives a $2^n$-cell partition of
$V^{\otimes n}$ by the Hamming weight of the sign-tuple.
\end{itemize}
We refer to the $2^n$ cells as **fine cells**; the cell at sign-weight
$k \in \{0, 1, \ldots, n\}$ has $\binom{n}{k}$ fine cells.

%==============================================================
\section{The dimensional ladder}
%==============================================================

\begin{theorem}[Dimensional ladder]\label{thm:ladder}
For $n = 0, 1, 2, 3, 4, 5$:
\[
\dim_{\F_5} V^{\otimes n} = 4^n = 2^{2n} = \dim_\R \Cl(2n).
\]
\end{theorem}

\begin{proof}
$\dim_{\F_5} V^{\otimes n} = (\dim_{\F_5} V)^n = 4^n$ by tensor-product
dimension. $4^n = (2^2)^n = 2^{2n}$. $\dim_\R \Cl(2n) = 2^{2n}$ by the
classical formula for the dimension of a Clifford algebra over a
quadratic form on $2n$ generators.
\end{proof}

The match is forced by the dimension of $V$. The non-trivial content is
the cell-decomposition correspondence:

\begin{theorem}[Binomial $\leftrightarrow$ grade correspondence]\label{thm:binomial}
For $n = 0, 1, \ldots, 5$, the $2^n$ fine cells of $V^{\otimes n}$
partition into groups of size $\binom{n}{k}$ for $k = 0, 1, \ldots, n$,
matching the grade-$k$ subspace dimensions of $\Cl(2n)$:
\[
\dim \Cl^{(k)}(2n) = \binom{2n}{k}, \quad k = 0, 1, \ldots, 2n.
\]
\end{theorem}

\begin{remark}
Note that $\binom{n}{k}$ (cells of $V^{\otimes n}$ at sign-weight $k$)
differs from $\binom{2n}{k}$ (grade-$k$ subspace of $\Cl(2n)$). The
relationship: each tensor factor of $V$ contributes 2 bits of "structural
sign" (idempotent class membership), giving $2n$ bits across $n$ factors,
so the sign-weight of a fine cell in $V^{\otimes n}$ corresponds to a
multi-grade subspace in $\Cl(2n)$ via the canonical embedding.
\end{remark}

%==============================================================
\section{Special case $n = 5$: the SU(5) representation match}
%==============================================================

\begin{theorem}\label{thm:su5}
At $n = 5$, $V^{\otimes 5}$ has 32 fine cells with the binomial
decomposition $1 + 5 + 10 + 10 + 5 + 1$. This decomposition matches the
representation content
\[
\mathbf{1} \oplus \bar{\mathbf{5}} \oplus \mathbf{10}
\quad \oplus \quad
\mathbf{1} \oplus \mathbf{5} \oplus \bar{\mathbf{10}}
\]
of SU(5), which is the representation content for one Standard Model
fermion generation plus its antimatter conjugate
(\cite{GeorgiGlashow1974}).
\end{theorem}

\begin{proof}
The 32-cell count and binomial decomposition are verified in the
accompanying script \texttt{test\_tig\_dirac.py}, test T14. The
correspondence to the SU(5) reps follows from
$\dim \mathbf{1} = 1 = \binom{5}{0}$, $\dim \bar{\mathbf{5}} = 5 = \binom{5}{1}$,
$\dim \mathbf{10} = 10 = \binom{5}{2}$, and the conjugation symmetry
$\binom{5}{k} = \binom{5}{5-k}$.
\end{proof}

This match is dimensional. We do not assert that $V^{\otimes 5}$ carries
the SU(5) action canonically; that question is left for future work.

%==============================================================
\section{Higher levels: status and conjecture}
%==============================================================

\begin{theorem}[Verified levels]
The dimensional ladder $\dim V^{\otimes n} = \dim_\R \Cl(2n)$ holds for
$n = 0, 1, 2, 3, 4, 5$, verified in test T15 of \texttt{test\_tig\_dirac.py}.
\end{theorem}

\begin{remark}[Conjectural extension]
We expect the dimensional ladder to hold for all $n \geq 0$ by induction
on the periodicity $\Cl(n+8) \cong \Cl(n) \otimes \Cl(8)$. Explicit
verification at $n = 6, 7, \ldots$ requires more compute than the
$<2$-second budget of the current verification scripts but is otherwise
straightforward.
\end{remark}

%==============================================================
\section{Verification}
%==============================================================

The accompanying scripts verify Theorems~\ref{thm:ladder},
\ref{thm:binomial}, and \ref{thm:su5}:
\begin{itemize}
\item \texttt{test\_tig\_dirac.py} test T15: dimensional ladder for
$n = 0, 1, 2, 3, 4, 5$.
\item Test T13: $V^{\otimes n}$ has $2^n$ fine cells for $n = 1, \ldots, 5$.
\item Test T14: $V^{\otimes 5}$ binomial $1+5+10+10+5+1 = 32$.
\end{itemize}
All three tests pass deterministically in $<1$ second on standard
hardware. The reference Python library \texttt{tig\_dirac.py} is
deposited at the public repository:
\url{https://github.com/TiredofSleep/ck/tree/tig-synthesis/Gen12/targets/clay/papers/sprint18_bridge_dirac_2026_05_04}.

%==============================================================
\section*{Acknowledgments}
%==============================================================

This work is part of a broader algebraic and physical framework developed
at 7Site LLC (Trinity Infinity Geometry, public repository at
github.com/TiredofSleep/ck). The present manuscript is self-contained as
a study in tensor algebras over finite fields; familiarity with the
broader framework is not required.

%==============================================================
\begin{thebibliography}{99}
%==============================================================

\bibitem{HesteneSobczyk1984} D.~Hestenes, G.~Sobczyk,
\emph{Clifford Algebra to Geometric Calculus},
D.~Reidel Publishing Co., 1984.

\bibitem{GeorgiGlashow1974} H.~Georgi, S.~L.~Glashow,
\emph{Unity of all elementary-particle forces},
Phys.~Rev.~Lett.~\textbf{32} (1974), 438--441.

\bibitem{Bott1959} R.~Bott,
\emph{The stable homotopy of the classical groups},
Ann.~Math.~\textbf{70} (1959), 313--337.

\bibitem{Sanders2026Dirac} B.~R.~Sanders,
\emph{A discrete Dirac-type algebra on $\F_5^4$: structural features,
field-invariance, and a Clifford-algebra dimensional ladder},
Manuscript in preparation, 2026.

\bibitem{TIGsynthesis2026} B.~R.~Sanders,
\emph{Trinity Infinity Geometry Synthesis},
github.com/TiredofSleep/ck (default branch),
DOI: 10.5281/zenodo.18852047, 2026.

\end{thebibliography}

\end{document}
